\section{Related work}


% \paragraph{Path-based Reasoning on Knowledge Graphs.} Path-based reasoning has emerged as a principal paradigm for interpretable multi-hop KGQA. Typically, this approach formulates the reasoning process as a sequential decision-making task, leveraging Reinforcement Learning (RL) to train agents to traverse the graph~\citep{xiong2017deeppath, das2018go}. Although foundational research has demonstrated the effectiveness of this methodology, challenges such as efficient exploration and the mitigation of sparse rewards still persist~\citep{jiang2023path, zhan2024reasoning}. To address these challenges, recent works have introduced more advanced search strategies, MCTS, which provides a more effective balance between exploration and exploitation~\citep{long2025enhancing}. However, a key limitation of these approaches is that the scoring function are pre-defined or fixed during inference, preventing adaptation to the evolving, query-specific context that arises throughout the MCTS expansion~\citep{shen2025reasoning}. Consequently, the search process may stray from optimal reasoning trajectories, ultimately compromising answer accuracy. In this paper, we propose an adaptive mechanism within MCTS that leverages paths generated during the MCTS expansion as contextual signals to fine-tune the path-scoring model, enabling continual adaptation to evolving contexts within each search episode.

% \paragraph{LLMs for KGQA.}
% Large Language Models (LLMs) have transformed KGQA by enabling flexible, interpretable multi-hop reasoning grounded in powerful language priors~\citep{pei2025mathfusion,zhao2023divknowqa}. Early approaches typically follow a retrieve-then-reason paradigm, extracting candidate paths using GNNs or heuristic methods~\citep{ma2025debate, fang2024karpa} and ranking them with LLMs~\citep{sui2024fidelis, li2024framework}. However, their reliance on static candidate sets limits adaptability to new questions and evolving contexts~\citep{sen2023knowledge, tian2024augmenting}. More recent strategies such as in-context learning and Chain-of-Thought prompting seek to unify retrieval and reasoning, but remain vulnerable to hallucinations and instability in domain-specific scenarios~\citep{huang2025survey, wang2024factuality}. MCTS-based frameworks further improve search efficiency and factual grounding by incrementally expanding and evaluating candidate paths~\citep{silver2016mastering, yao2023tree}, yet typically use static evaluation models that cannot adapt to dynamic semantic contexts~\citep{yuan2024advancing, bi2024forest}. In contrast to these approaches, this work investigates the development of an adaptive MCTS-based framework for KGQA, in which the evaluation model is continuously updated online to enable robust, context-aware reasoning across diverse domains.

\paragraph{Knowledge Graph Question Answering (KGQA).} KGQA aims to enhance reasoning capabilities by incorporating external knowledge graphs to answer natural language questions~\citep{choi2023nutrea, xu2025memory}. Existing KGQA approaches can be broadly classified into two categories: retrieve-then-reason and dynamic path generation. The first category extracts candidate reasoning paths using Graph Neural Networks (GNNs)\citep{ma2025debate, yao2025learning} or rule-based heuristics\citep{fang2024karpa}, followed by LLM-based answer generation. While GNNs learn embeddings to identify relevant paths and rule-based methods apply predefined patterns~\citep{zhao2023divknowqa, liu2025dual}, these approaches lack the flexibility to adapt dynamically to question-specific context during inference. 
In contrast, dynamic path generation methods, such as CoT prompting~\citep{sui2024fidelis, li2024framework} and MCTS~\citep{ma2025deliberation, shen2025reasoning}, unify retrieval and reasoning for more flexible exploration. However, they suffer from high computational overhead due to repeated LLM calls, and static scorers often fail to adapt to evolving path semantics~\citep{long2025enhancing, luo2025kbqa,ma2024think}. To address these challenges, we propose an adaptive framework that integrates symbolic search with a fine-tuned evaluation model, aiming to improve both computational efficiency and reasoning accuracy in KGQA.


\noindent\textbf{Adaptive and Self-Improving Reasoning Models.} A promising approach to developing adaptive reasoning models is to frame the process within a reinforcement learning (RL) paradigm, where an agent learns a policy to navigate a state space. Early methods such as DeepPath~\citep{xiong2017deeppath} and MINERVA~\citep{das2018go} used RL to discover reasoning paths by rewarding the agent only when a correct answer is reached. However, this leads to the sparse rewards problem, as positive feedback arrives only after long action sequences, resulting in weak learning signals and poor exploration efficiency~\citep{zhai2024fine,chang2023learning}. To address this challenge, an alternative is self-training via pseudo-labeling, where the model learns from its own high-confidence predictions~\citep{lee2013pseudo,xie2020self}. While commonly used in semi-supervised learning, pseudo-labeling proves especially effective in reasoning tasks with limited supervision~\citep{wang2022semi,huang2025mllm}. Instead of relying on sparse terminal rewards, we leverage intermediate search paths as dynamic pseudo-paths, offering dense and adaptive supervision. This facilitates continual refinement of the path evaluator to better capture the evolving semantics of reasoning.